\documentclass[11pt,a4paper]{article}

\usepackage[T1]{fontenc}
\usepackage[utf8]{inputenc}
\usepackage{lmodern}
\usepackage[margin=2.7cm]{geometry}
\usepackage{amsmath}
\usepackage{amssymb}
\usepackage{amsthm}
\usepackage{booktabs}
\usepackage{enumitem}
\usepackage{algorithm}
\usepackage{algpseudocode}
\usepackage{float}
\usepackage{microtype}
\usepackage[hidelinks]{hyperref}

\newtheorem{definition}{Definition}[section]
\newtheorem{proposition}[definition]{Proposition}
\newtheorem{lemma}[definition]{Lemma}
\newtheorem{remark}[definition]{Remark}
\newtheorem{example}[definition]{Example}

\newcommand{\N}{\mathbb{N}}
\newcommand{\Npos}{\mathbb{N}_{>0}}
\newcommand{\Pow}{\mathcal{P}}
\newcommand{\pre}{\operatorname{pre}}
\newcommand{\post}{\operatorname{post}}
\newcommand{\inh}{\operatorname{inh}}
\newcommand{\Enabled}{\operatorname{enabled}}
\newcommand{\Succ}{\operatorname{Succ}}
\newcommand{\Pred}{\operatorname{Pred}}
\newcommand{\Reach}{\operatorname{Reach}}
\newcommand{\Seeds}{\Sigma}
\newcommand{\Off}{R}
\newcommand{\kind}{\operatorname{kind}}
\newcommand{\cls}{\operatorname{class}}
\newcommand{\hits}{\operatorname{hits}}
\newcommand{\need}{\operatorname{need}}
\newcommand{\intentional}{\operatorname{intentional}}
\newcommand{\vis}{\operatorname{vis}}
\newcommand{\step}{\operatorname{step}}
\newcommand{\mix}{\operatorname{mix}}
\newcommand{\FNV}{\operatorname{FNV}}
\newcommand{\TS}{\mathcal{H}}
\newcommand{\Pair}{\mathcal{P}}
\newcommand{\Tau}{\tau_{\max}}
\newcommand{\wgt}{\operatorname{w}}
\newcommand{\dist}{\operatorname{dist}}
\newcommand{\tierpen}{\operatorname{penalty}}
\newcommand{\R}{\mathbb{R}}
\newcommand{\B}{\mathbb{B}_{32}}
\newcommand{\fir}[1]{\xrightarrow{#1}}
\newcommand{\bool}{\{0,1\}}

\title{\textbf{Ant Colony Optimization for Coverage in Model-Based Testing}}
\author{Esa Kettunen \\
  Egge Collective \\
  The \texttt{eggembt} project}
\date{Draft of 25 November 2026 --- product version \texttt{v0.0.8}}

\begin{document}

\maketitle

\begin{abstract}

We give a mathematical account of the modelling language and
exploration heuristic used by \texttt{eggembt}, a Petri-net tool for
model-based testing, optimized for model coverage.  We define the
class of nets the tool accepts (basically places with bounds, weighted
arcs, inhibitor arcs), augmented with types and labels for transitions,
several named initial markings, as well as target and error markings,
milestones, rare transitions, and crash predictions. We also define
the lazy reachability graph---the state space of the net---which the
tool builds during testing, the coverage and target machinery that
turns the graph into a testing plan, and the ant-colony optimisation
(ACO) that steers the plan.  The heuristic is presented as a
stochastic process on the graph: two learned pheromone channels, one
static \emph{off-target} channel, novelty weighting, back-propagating
deposits, evaporation, and an optional finer store indexed by a
marking together with a hash of the recent trace (a \emph{T-suffix}).
The final section relates the suggester's score to the colony.  The
definitions are complemented by pseudo-code for graph generation, ant
movement, pheromone updates and event selection.  Every definition is
taken from the implementation, not invented for the paper; the default
parameters quoted here are the ones in the source at version
\texttt{v0.0.8}.

The tool is intended merely as a proof of concept but it also has some
characteristics that show how model-based testing can be integrated into
the software process of a real-world organization.

The tool and its documentation were created with the help of
artificial intelligence. Curiously enough, the original idea,
prompting, nudging, hand-waving, and overall project oversight
(including but not limited to this very sentence) were the
responsibility of the author, a human, and so will be any errors to be
found herein.

\end{abstract}

\medskip
\noindent\textbf{Keywords:} Petri nets; model-based testing; reachability
graph; state space; coverage; ant colony optimization; test generation.

\tableofcontents

\section{Introduction}

\subsection{Scope}

\texttt{eggembt} tests a system under test (SUT) against a Petri-net
specification.  Petri nets, introduced by Petri
\cite{petri1962kommunikation}, are a classical model of concurrent and
distributed systems \cite{murata1989petri}, and deriving tests from a
model is the subject of model-based testing
\cite{utting2007practical,utting2012taxonomy}.  The net supplies an
oracle: a step of the SUT that the net cannot explain is either a bug
candidate or a specification error, and the part of the net that no
run has visited is untested.  The tool explores the reachability graph
(RG) of the net---the graph of every marking the net can reach, also
called its state space---while it drives the SUT, records coverage,
and suggests the next input.  The suggestion (which can be considered
as the selection of the next test step) is the place where the ant
colony enters: virtual ants walk the \emph{model}, never the SUT, and
leave traces of pheromone making productive or promising routes more
likely in later walks.  The heuristic is an ant-colony optimisation
\cite{dorigo1996ants,dorigo2004aco} over the reachability graph.

This document formalises exactly that pipeline.  It is deliberately close to
the code.  Where the implementation makes a choice that a textbook definition
would leave open---for example what happens when a firing would exceed a place
bound, or how a target marking differs from a target transition---the choice
is stated as part of the definition.  Sections~\ref{sec:nets} and
\ref{sec:rg} define the net and the graph, Sections~\ref{sec:targets}
and~\ref{sec:offtarget} define targets and the off-target region,
Section~\ref{sec:aco} defines the colony, Section~\ref{sec:suggester} defines
the suggester's score, Section~\ref{sec:algorithms} restates the procedures as
pseudo-code, Section~\ref{sec:application} revisits the campaigns recorded in
the journal, and Section~\ref{sec:implementation} maps symbols to the
implementation.
The companion design record, \texttt{journal.md}, contains the engineering
history and the reasons behind the choices.

\subsection{Notation}

For a set $X$, $\Pow(X)$ is its power set and $\Pow_{\mathrm{fin}}(X)$ the set
of finite subsets.  A \emph{multiset} over $X$ is a function $X\to\N$ with
finite support.  We write $M(p)$ for the multiplicity of $p$ and identify a
multiset with its graph.  $\mathbb{1}_{A}$ is the indicator of a predicate
$A$.  All token counts are 32-bit unsigned integers in the implementation;
the mathematical definitions use $\N$ and note the machine limit where it
matters.  We write $\B=\{0,\dots,2^{32}-1\}$ for the 32-bit words.  A partial
function $f\colon X\rightharpoonup Y$ has a domain
$\operatorname{dom}(f)$.  We mostly say \emph{marking} for an element of the
state space; the word \emph{state} is kept for the compound \emph{state space},
a synonym for the reachability graph, and for implementation names such as
\texttt{maxStates}.

\section{Petri nets with tester semantics}\label{sec:nets}

Our nets are Place/Transition nets with weighted ordinary arcs, weighted
inhibitor arcs, place bounds and typed transitions, in the classical sense of
\cite{petri1962kommunikation,murata1989petri}; the annotations of
Section~\ref{sec:annotations} are specific to the testing tool.

\subsection{The net}

\begin{definition}[Net]\label{def:net}
A \emph{net} is a tuple
\[
  N=(P,T,F,I,w,\beta,\ell,\kappa)
\]
where
\begin{enumerate}[label=(\roman*),leftmargin=2.2em]
  \item $P$ and $T$ are finite, non-empty, disjoint sets of \emph{places} and
        \emph{transitions};
  \item $F\subseteq (P\times T)\cup(T\times P)$ is the set of \emph{ordinary
        arcs};
  \item $I\subseteq P\times T$ is the set of \emph{inhibitor arcs};
  \item $w\colon F\cup I\to\Npos$ assigns a positive \emph{weight} to every
        arc;
  \item $\beta\colon P\to\Npos\cup\{\infty\}$ assigns a \emph{bound} to every
        place, with $\infty$ meaning unbounded;
  \item $\ell\colon T\to\Sigma^{*}$ assigns a label over an alphabet $\Sigma$;
        $\ell(t)=\varepsilon$ means that $t$ is \emph{unlabelled};
  \item $\kappa\colon T\to\{\mathsf{C},\mathsf{O},\mathsf{I}\}$ assigns a
        \emph{kind}: $\mathsf{C}$ controllable (the tester may send it),
        $\mathsf{O}$ observable (the SUT emits it), $\mathsf{I}$ internal.
\end{enumerate}
A place may have both an ordinary and an inhibitor arc to the same transition;
the two conditions are then imposed independently.
\end{definition}

The label of a transition is the SUT-level event it corresponds to.  An
unlabelled internal transition is \emph{silent}: it changes the marking but
cannot be reported by a plug-in.  An internal transition with a label is
\emph{traced} by default, and a model may explicitly override visibility;
write $\vis(t)$ for the resulting visibility predicate.  Observable
transitions are always visible; silent internal transitions are never visible
to the tester.

For a transition $t$ and place $p$ define the three arc multiplicities
\[
  \pre_t(p)=
  \begin{cases}
    w(p,t), & (p,t)\in F,\\
    0, & \text{otherwise},
  \end{cases}
  \qquad
  \post_t(p)=
  \begin{cases}
    w(t,p), & (t,p)\in F,\\
    0, & \text{otherwise},
  \end{cases}
\]
and
\[
  \inh_t(p)=
  \begin{cases}
    w(p,t), & (p,t)\in I,\\
    \infty, & \text{otherwise}.
  \end{cases}
\]
The use of $\infty$ makes the inhibitor test vacuously true when no inhibitor
arc exists.  A \emph{marking} is a multiset $M\colon P\to\N$.  The
multiplicity $M(p)$ is the number of tokens on $p$.

\begin{definition}[Enabledness]\label{def:enabled}
Transition $t$ is \emph{enabled} in marking $M$, written $M\models
\Enabled(t)$, iff
\[
  \forall p\in P:\quad
  M(p)\ge \pre_t(p)
  \quad\text{and}\quad
  M(p)<\inh_t(p).
\]
\end{definition}

Enabledness deliberately ignores place bounds.  A transition may therefore be
enabled while its firing would violate a bound; such a firing is a
\emph{bound violation}, not a step (Definition~\ref{def:firing}).

\begin{definition}[Firing]\label{def:firing}
Let $M\models\Enabled(t)$ and define the candidate marking
\[
  M' = M-\pre_t+\post_t,
  \qquad
  M'(p)=M(p)-\pre_t(p)+\post_t(p)\quad\text{for all }p\in P.
\]
The \emph{firing step} $M\fir{t}M'$ exists iff
\[
  M'(p)\le \beta(p)\quad\text{for every }p\in P
\]
and every $M'(p)$ fits into the implementation's 32-bit token counter.  If
the bound or the counter is exceeded, the attempted firing is recorded as a
\emph{bound violation} $(M,t,p)$ and creates no step.  Tokens are consumed
before they are produced, so a place that is both an input and an output of
$t$ is handled by the single formula above.
\end{definition}

The step relation is the usual Petri-net firing relation with one tester
convention: a firing that would exceed a declared bound is not silently
clamped and not turned into a transition to an error marking; it is excluded
from the reachability graph and reported separately.  This keeps a cut-off
model honest: the graph never invents a marking the model forbids.

\subsection{Named markings and classification}

The model distinguishes three roles for named markings.

\begin{definition}[Named markings]\label{def:named}
An \emph{initial marking} is a marking $M_0$.  A \emph{target marking} is a
pair $(A,\mu)$ with $A\in\Pow_{\mathrm{fin}}(P\times\Npos)$ and a match mode
$\mu\in\{\mathsf{exact},\mathsf{covers}\}$.  An \emph{error marking} is a
target marking equipped with a policy
$\pi\in\{\mathsf{until\_covered},\mathsf{once},\mathsf{avoid}\}$, an optional
hit count, a dump mode and a dump budget.  The model supplies finite families
\[
  \mathcal{M}_0=\{M_0^{(1)},\dots,M_0^{(k)}\},\quad
  \mathcal{G}=\{(A_j,\mu_j,\iota_j)\}_{j=1}^{m},\quad
  \mathcal{E}=\{(A_i,\mu_i,\pi_i,r_i)\}_{i=1}^{n},
\]
where $\iota_j\in\bool$ marks a target marking as \emph{intentional}.
\end{definition}

A finite marking $A$ is read as the marking that puts $A(p)$ tokens on each
listed place and zero tokens everywhere else.  The match relation is
\[
  M\models_{\mathsf{exact}} A \iff M=A,
  \qquad
  M\models_{\mathsf{covers}} A \iff \forall p\in P:\; M(p)\ge A(p).
\]
Error markings are tested before target markings.  The classification of a
marking is
\[
  \kind(M)=
  \begin{cases}
    \mathsf{Error}(i), & \text{the least }i\text{ with }M\models_{\mu_i} A_i,\\
    \mathsf{Target}(j), & \text{the least }j\text{ with }M\models_{\mu_j} A_j
                         \text{ and no error matches},\\
    \mathsf{Normal}, & \text{otherwise}.
  \end{cases}
\]
The least-index convention is the implementation's: when several named
markings match, the first in file order wins, and an error always wins over a
target.  A marking with kind $\mathsf{Target}$ or $\mathsf{Error}$ is never
reported as a deadlock merely because no transition is enabled; it is a
declared end or fault.

\subsection{Transition annotations}\label{sec:annotations}

The net also carries annotations that do not change the token game but change
how the tester interprets a step.  They are part of ``our kind of Petri net''
because the exploration heuristics read them directly.

\begin{center}
\small
\begin{tabular}{@{}p{0.24\linewidth}p{0.68\linewidth}@{}}
\toprule
Annotation & Meaning \\
\midrule
$\mathsf{error}(t)$ &
The transition itself is an error.  A \emph{certain} step of $t$ fails the
run, requests a post-mortem, and lays catastrophe pheromone. \\
$\mathsf{milestone}(t)=r$ &
A certain step of $t$ is a success of rank $r\in\Npos$, worth an attractive
deposit of $1.5\,r$.  It is not the campaign's goal, only a signpost. \\
$\mathsf{rare}(t)=(p,a,\delta)$ &
The transition is declared to fire with small probability $p$, may be hunted
for at most $a$ attempts, and a dwell at its enabling marking lasts $\delta$
milliseconds.  The probability is a promise by the modeller, not measured. \\
$\mathsf{crash}(t)=c$ &
The SUT may die while stepping $t$.  The crash class $c$ (always, often,
sometimes, rare, random), an optional prior mean and a reproduction budget
calibrate the credibility experiment that turns a probable crash into a
confirmed one. \\
$\mathsf{command}(t)$ &
A shell command executed when a step of $t$ is certain; an optional place
guard and an execution budget are part of the annotation. \\
$\mathsf{lifecycle}(t)$ &
A controllable transition whose label is one of the reserved words
\textsf{start}, \textsf{restart}, \textsf{stop}, \textsf{pause},
\textsf{resume}.  The tester still sends the label; with no output arcs the
engine, not the token game, supplies the landing marking.  Such a transition
counts as enabled but creates no ordinary RG edge. \\
\bottomrule
\end{tabular}
\end{center}

\noindent
A \emph{certain} step is one on which the belief tracker can name the fired
transition, not merely one of several candidates; Section~\ref{sec:belief}
recalls the belief set.  Crash predictions may also be attached to a named
marking, in which case the markings matching it play the role of the
predicted place.  A predicted crash is a \emph{target} until the reproduction
experiment resolves it; only then does the catastrophe channel take over.

\subsection{Belief sets}\label{sec:belief}

A run does not observe the marking directly.  It sees a sequence of visible
events: stimuli it sent, outputs and traced internal steps reported by the
plug-in, and quiescence reports.  The tracker therefore maintains a
\emph{belief set} $B\subseteq\Reach(N)$ of markings consistent with the events
so far.  Initially $B$ is the set of chosen initial markings (or a subset if
the campaign restricts the initials).  A stimulus with label $a$ followed by
an observed label $b$ updates $B$ by the relation
\[
  B' = \bigl\{\,M'\;:\;\exists M\in B,\; t,t'\in T,\;
        \ell(t)=a,\ \ell(t')=b,\;
        M\fir{t}M''\fir{t'}M'\,\bigr\},
\]
with the obvious variants for silent internal steps and for an observed event
that needs no preceding stimulus.  When an observed label is possible through
several transitions, all of them remain candidates.  The implementation also
records \emph{ambiguous} coverage for elements that lie on some but not all
candidate paths.

The belief matters to the colony in two ways.  First, the ants start from the
belief markings, so their walks explore the part of the model the SUT may
currently occupy.  Second, deposits from real runs are anchored at definitely
credited steps: a milestone or an error transition deposits only when the
belief pins the step down.

\section{Reachability graph (state space), coverage and targets}\label{sec:rg}

\subsection{The full reachability graph}

\begin{definition}[Reachability graph]\label{def:rg}
The \emph{reachability graph}---also called the \emph{state space}---of
$N$ from a set $\mathcal{M}_0$ of initial markings is the labelled directed
graph
\[
  G(N,\mathcal{M}_0)=(S,E,S_0),
\]
where $S\subseteq\N^{P}$ is the least set containing every $M_0^{(i)}$ and
closed under firing, $S_0=\{M_0^{(i)}\}$ is the set of chosen initial
markings, and
\[
  E=\{\,(M,t,M')\;:\;M\in S,\ M\fir{t}M'\,\}.
\]
Nodes are distinct markings, not occurrences: two different paths that reach
the same marking share one node.  The reachability graph is therefore exactly
the state space of the net.  An edge is written $e=(M,t,M')$ with source
$\sigma(e)=M$, target $\tau(e)=M'$ and label $\ell(t)$.
\end{definition}

The graph is generally infinite even for a finite net, and even a finite
graph can be far too large to enumerate; this is the state explosion problem
\cite{valmari1998state}.  The tool therefore never constructs $G$ in full.
It constructs a finite \emph{explored subgraph} and treats the unexplored
part explicitly.

\begin{definition}[Explored subgraph]\label{def:explored}
An explored subgraph is a triple $(S_{\mathrm{known}},S_{\mathrm{exp}},E_{\mathrm{exp}})$
such that
\begin{enumerate}[label=(\roman*),leftmargin=2.2em]
  \item $S_{\mathrm{exp}}\subseteq S_{\mathrm{known}}\subseteq S$;
  \item every edge of $E_{\mathrm{exp}}$ has its source in $S_{\mathrm{exp}}$
        and its target in $S_{\mathrm{known}}$;
  \item for every $s\in S_{\mathrm{exp}}$, every legal successor of $s$ is in
        $E_{\mathrm{exp}}$, unless the marking limit stopped the creation of
        a new node, in which case $s$ is flagged \emph{incomplete}.
\end{enumerate}
The \emph{frontier} is $U=S_{\mathrm{known}}\setminus S_{\mathrm{exp}}$.
\end{definition}

The marking limit (the implementation's \texttt{maxStates}, default one
million markings) is a hard cap.  When it is hit, new nodes are not created,
the source marking is flagged incomplete, and the graph is marked truncated.
Analyses must never confuse the resulting frontier with a deadlock.  Bound
violations are recorded beside the graph and never become edges.

\subsection{Coverage and saturation}

Coverage is accumulated over runs and restored from the run log.  For each
marking, edge and transition the database keeps a \emph{definite} count (the
visit lies on the one known path) and an \emph{ambiguous} count (the visit is
possible, but the belief does not pin it down).  Write
\[
  d(s),\ d(e),\ d(t)\in\N
\]
for the definite counts and $\theta\in\Npos$ for the coverage threshold
(default $\theta=1$).

\begin{definition}[Saturation and novelty]\label{def:saturation}
An element is \emph{saturated} when its definite count reaches $\theta$.  For
an edge $e$ the \emph{novelty} is
\[
  \nu(e)=
  \begin{cases}
    \dfrac{1}{1+d(e)}, & d(e)<\theta,\\[2mm]
    0.01, & d(e)\ge\theta.
  \end{cases}
\]
\end{definition}

Novelty is the colony's ``unexplored smells interesting'' term.  The value
$0.01$ is deliberately tiny but nonzero: a saturated edge is not absolutely
forbidden, it merely stops competing with unexplored ones.  At the coverage
level, saturation also wipes an edge's attractive pheromone during
evaporation.

\subsection{Target phases and board classes}\label{sec:targets}

The target board turns the static target and error markings into a dynamic
classification of markings.

\begin{definition}[Target phase]\label{def:phase}
For a target or error marking with required definite visits
$r\in\N$ and observed definite visits $h\in\N$, the phase is
\[
  \operatorname{phase}=
  \begin{cases}
    \mathsf{Seeking}, & h<r,\\
    \mathsf{Satisfied}, & h\ge r.
  \end{cases}
\]
\end{definition}

For an error marking the requirement follows its policy:
\[
  r=
  \begin{cases}
    0, & \pi=\mathsf{avoid},\\
    1, & \pi=\mathsf{once},\\
    r_{\text{declared}}\ \text{or}\ \theta, & \pi=\mathsf{until\_covered}.
  \end{cases}
\]
Target markings are intentional or not.  An intentional target marking has
$r=1$; the first certain visit ends the campaign.  A non-intentional target
marking is not a goal and is never forbidden; it remains an idle or
``legitimate end'' marking used by the livelock analysis.

Deadlocks and certain terminal cycles are additional targets.  All deadlocks
of the explored graph form one target, with the policy configured for
deadlocks (default \textsf{once}).  Every certain terminal cycle found by the
static livelock analysis is one target, again with its own policy and hit
count.  While a crash prediction is unresolved, the markings where the
predicted transition is enabled (or the markings that certainly match the
predicted named marking) are \emph{seeking}, overriding an earlier
classification.

\begin{definition}[Board class]\label{def:class}
The \emph{board class} of a marking $s$ is
\[
  \cls(s)=
  \begin{cases}
    \mathsf{Forbidden}, & \text{satisfied target or \textsf{avoid} error},\\
    \mathsf{Seeking}, & \text{still-sought target, error, deadlock,
                        livelock or crash},\\
    \mathsf{Normal}, & \text{otherwise}.
  \end{cases}
\]
Error markings take precedence over target markings; a marking that belongs
to both inherits the error's class.
\end{definition}

\begin{remark}
``Forbidden'' is a steering decision, not a claim that the marking is
illegal.  A satisfied error marking has been seen often enough; the campaign
now avoids
the expensive post-mortem and restart.  A deadlock or terminal cycle that has
been confirmed is treated the same way.  The run outcome still reports the
situation when it happens; only the planner's preference changes.
\end{remark}

\subsection{The off-target region}\label{sec:offtarget}

The target markings describe where a run should end.  Their complement---the
markings from which no such end is reachable any more---is a
\emph{reward-inaccessible} region of the explored graph: a run can enter it
but never legitimately leave.  The tool calls it the \emph{off-target region}
and marks its entrances with a static repellent value.

\begin{definition}[Off-target region]\label{def:offtarget}
Let $G_{\mathrm{exp}}$ be the explored subgraph and let
\[
  \begin{aligned}
    \Seeds={}&
    \{\,s\in S_{\mathrm{exp}}:\kind(s)=\mathsf{Target}\,\}\\
    &{}\cup
    \{\,s\in S_{\mathrm{exp}}:\cls(s)=\mathsf{Seeking}\,\}\\
    &{}\cup
    \{\,s\in S_{\mathrm{exp}}:
        s\text{ is incomplete or has an edge into }U\,\}.
  \end{aligned}
\]
Define the \emph{off-target region} as
\[
  \Off=\{\,s\in S_{\mathrm{exp}}:\text{there is no path in }
        G_{\mathrm{exp}}\text{ from }s\text{ to a marking in }\Seeds\,\}.
\]
Equivalently, $\Off$ is the complement in $S_{\mathrm{exp}}$ of the set of
markings that can reach $\Seeds$ by a (possibly empty) path.  It is computed by
a reverse breadth-first search from $\Seeds$ over the predecessor index of the
explored graph.
\end{definition}

The three kinds of seeds are the markings behind which a legitimate end is
\emph{known} to exist or the future is \emph{unknown}:

\begin{itemize}
  \item every target-kind marking, including a non-intentional one, because it
        is a legitimate end for the livelock analysis;
  \item every marking the board is still seeking---an error marking, a
        deadlock, a livelock, an unresolved crash prediction---because the
        campaign goes there on purpose;
  \item every expanded marking that can still step into the frontier or whose
        expansion was cut off, because a target marking might lie beyond what
        is not known yet.
\end{itemize}

\begin{definition}[Off-target channel]\label{def:offtarget-channel}
Let $\rho\ge0$ be the configured off-target repel (default $\rho=1.0$;
$\rho=0$ switches the mechanism off).  The static channel
$g\colon E_{\mathrm{exp}}\to[0,\rho]$ is
\[
  g(e)=
  \begin{cases}
    \rho, & \tau(e)\in\Off,\\
    0, & \text{otherwise}.
  \end{cases}
\]
\end{definition}

The channel is \emph{set}, never accumulated and never evaporated.  It is
recomputed from scratch whenever the graph has grown, since a newly expanded
marking can turn a previously lost marking into one that can still reach a
target.  Edges into the frontier are deliberately not marked: unknown
territory is presumed innocent until it is expanded.

\begin{proposition}[The region is closed under successors]\label{prop:closed}
Let $s\in\Off$ and let $(s,t,s')\in E_{\mathrm{exp}}$.  Then $s'\in\Off$.
Consequently a walk that enters $\Off$ can never leave it through an explored
edge, and it can never reach a marking of $\Seeds$ or the frontier.
\end{proposition}

\begin{proof}
If $s'$ could reach a seed, then $s$ could reach it by the edge $(s,t,s')$,
contradicting $s\in\Off$.  If $s'$ were in the frontier, then $s$ would have
an edge into $U$ and would itself be a seed, again a contradiction.  Hence
$s'\in S_{\mathrm{exp}}\setminus\Seeds$ and $s'$ cannot reach $\Seeds$; by
Definition~\ref{def:offtarget}, $s'\in\Off$.
\end{proof}

Thus $\Off$ is exactly the largest forward-closed set of expanded markings
that contains no rewarding seed.  The ants are not killed when they enter
it---uncovered behaviour inside the region may still be worth covering---but
every edge into the region carries the static repel $g(e)$, which lowers the
weight of that edge by the factor $e^{-\lambda\rho}$ in the choice rule
(Section~\ref{sec:choice}).  With the default $\lambda=1$ and $\rho=1$, the
factor is $e^{-1}\approx0.368$.

\section{The ant colony}\label{sec:aco}

\subsection{Channels and the environment}

Fix an explored graph $G_{\mathrm{exp}}=(S,E)$.  The colony keeps two learned
channels and the static off-target channel:
\[
  a\colon E\to[0,\Tau],\qquad
  r\colon E\to[0,\Tau],\qquad
  g\colon E\to[0,\rho],
\]
called \emph{attract}, \emph{repel} and \emph{off-target}.  Initially
$a(e)=r(e)=0$; $g$ is computed by Definition~\ref{def:offtarget-channel}.
The cap is $\Tau=10$ by default.  The channels are not persisted: edge
identifiers belong to one session's graph, and the colony re-learns quickly
from the coverage that the run log restores.

At a marking $s$, the outgoing edges split into the tester's inputs and the
SUT's own moves:
\[
  E_{\mathrm{C}}(s)=\{e\in E:\sigma(e)=s,\ \kappa(t_e)=\mathsf{C}\},
  \qquad
  E_{\mathrm{O}}(s)=\{e\in E:\sigma(e)=s,\ \kappa(t_e)\neq\mathsf{C}\}.
\]
Only controllable edges are under the tester's control; observable and
internal edges are the environment's moves.  With probability $p_{\mathrm{env}}$
(default $0.5$) an ant takes an environment move when both pools are nonempty;
if one pool is empty it takes the other.  Environment moves use a separate
weight $\omega(e)$, uniform by default and optionally the observed
offered-versus-taken rate of the transition's label.  The pheromone formula
does not apply to environment moves.

\subsection{The choice rule}\label{sec:choice}

\begin{definition}[Edge weight]\label{def:weight}
For an input edge $e=(s,t,s')$, let $\TS$ be the ant's current T-suffix
(Section~\ref{sec:tsuffix}) and let
\[
  \TS'=\TS\cdot h(t)
\]
be the suffix extended by the edge's step.  Write
$a^{+}(e)=\Pair_a(s',\TS')$ and $r^{+}(e)=\Pair_r(s',\TS')$ for the pair
channels (both zero when the T-suffix layer is off).  The \emph{weight} of the
input edge is
\[
  \wgt(e)=\bigl(a(e)+a^{+}(e)+0.1\bigr)^{\alpha}\,
        \nu(e)^{\beta}\,
        \exp\!\Bigl(-\lambda\bigl(r(e)+g(e)+r^{+}(e)\bigr)\Bigr),
  \label{eq:weight}
\]
with defaults $\alpha=\beta=\lambda=1$.  The constant $0.1$ is a floor: an
edge with no pheromone at all still has a positive base.
\end{definition}

Within the chosen pool, the probability of taking edge $e$ is proportional to
its weight:
\[
  \Pr(e\mid s)=
  \begin{cases}
    \dfrac{\wgt(e)}{\sum_{e'\in E_{\mathrm{C}}(s)}\wgt(e')},
      & e\in E_{\mathrm{C}}(s),\ \sum\wgt>0,\\[3mm]
    \dfrac{1}{|E_{\mathrm{C}}(s)|}, & e\in E_{\mathrm{C}}(s),\ \sum\wgt=0,\\[3mm]
    \dfrac{\omega(e)}{\sum_{e'\in E_{\mathrm{O}}(s)}\omega(e')},
      & e\in E_{\mathrm{O}}(s),
  \end{cases}
\]
with the environment case made uniform when the $\omega$-sum is zero.  The
two-step draw---environment versus input, then within the pool---is the exact
implementation; it keeps the SUT's own nondeterminism separate from the
tester's choice.

\subsection{Arrival, deposits and back-propagation}\label{sec:deposits}

An ant starts at a start marking, walks at most $40$ steps, and stops when
it arrives at a marking that teaches it something.  The arrival rule is:

\begin{enumerate}[label=(\roman*),leftmargin=2.2em]
  \item If $\cls(s)=\mathsf{Forbidden}$ (a catastrophe), the ant deposits the
        catastrophe amount $3.0$ into the repel channel and dies.
  \item If $\cls(s)=\mathsf{Seeking}$ (a still-sought target), the ant
        deposits the target amount $2.0$ into the attract channel and stops.
  \item If $s$ is unexpanded and the round's expansion budget is exhausted,
        the ant deposits the unexplored amount $1.0$ into the attract channel
        and stops.  If the budget is not exhausted, the marking is expanded
        and the walk continues.
  \item If $s$ has an outgoing edge whose target is not forbidden and whose
        definite coverage is zero, the ant deposits $1.0$ into the attract
        channel and stops: there is unexplored behaviour here.
  \item Otherwise the ant walks on.
\end{enumerate}

A deposit is not laid only on the ant's arrival edge.  From the arrival
marking $s$, the colony walks backwards through the predecessor index to a depth
$D\ge5$ (default $D=5$).  For an edge $e$ at reverse distance $d(e)$---where
$d(e)=0$ for the edges directly entering $s$, $d(e)=1$ for the edges entering
those predecessors, and so on---the deposit is
\[
  c(e)\leftarrow
  \min\!\Bigl(\Tau,\; c(e)+A\,\gamma^{d(e)}\Bigr),
  \qquad d(e)<D,
  \label{eq:backprop}
\]
where $A$ is the deposited amount, $\gamma=0.7$ is the depth decay, and
$c\in\{a,r\}$ is the channel.  Nodes are visited once, at their shortest
reverse distance, so a diamond in the graph is not double-counted.  In a
cycle the signal also reaches edges that merely lead back to the place where
the event can happen, weaker the further they are; the direct way in always
carries the most.

Real runs deposit the same channels with their own knowledge:
\[
  A_{\mathrm{milestone}}=1.5\,r,\qquad
  A_{\mathrm{catastrophe}}=3.0\bigl(1+0.25\min(k,8)\bigr),
\]
where $r$ is the milestone rank and $k$ the number of stimuli the run spent
on the way in (capped at eight).  A target marking that the campaign reached
deposits $2.0$.  A milestone is not priced by the way in---a success is a
success---while a catastrophe is: a crash one input away is marked lightly, a
crash deep in the model is marked hard.  The cap $\Tau$ remains the ceiling.
Deposits from real runs are anchored at definitely credited steps, so an
ambiguous stretch does not invent a milestone outcome.

\subsection{Evaporation, saturation and the matured-target exception}

After an iteration, when evaporation is on,
\[
  a(e)\leftarrow(1-\varepsilon_a)\,a(e),
  \qquad
  r(e)\leftarrow(1-\varepsilon_r)\,r(e),
  \label{eq:evaporation}
\]
with $\varepsilon_a=0.10$ and $\varepsilon_r=0.02$.  Repellent memory
therefore outlives attractive memory, which is deliberate: the lesson of a
costly mistake should last longer than the attraction to a goal that has
already been reached.  The static channel $g$ does not evaporate.

Saturation is a coverage rule, not evaporation: every saturated edge has its
attraction set to zero even if the rest is evaporating.  This is how the
colony moves on from a covered part of the graph.

The optional merge-time rule used when restoring pheromone from an earlier
campaign spares a matured target.  More precisely, every marking whose board
class is $\mathsf{Forbidden}$---a satisfied target, error, deadlock or
livelock---has its incoming edges (and the corresponding pair entries)
exempted from the per-iteration multiplication, while saturation still
applies.  The finding has reached its visited-$N$ criterion and the lesson
should outlive the merge.  The rule is off by default and is
orthogonal to the definitions above; it only changes which entries
Equation~\eqref{eq:evaporation} touches.

\subsection{The T-suffix layer}\label{sec:tsuffix}

The edge channels say ``this way is good''.  They do not distinguish two
routes that reach the same marking.  The optional T-suffix layer adds that
distinction with a bounded memory of the recent trace.

\begin{definition}[Step hash and suffix]\label{def:tsuffix}
For a transition $t$, the \emph{step identity} is
\[
  \operatorname{id}(t)=
  \begin{cases}
    \ell(t), & \ell(t)\neq\varepsilon,\\
    \text{``I:''}\,\Vert\,\operatorname{id}_{\mathrm{net}}(t), & \text{otherwise},
  \end{cases}
\]
and the \emph{step hash} is the 32-bit FNV-1a hash
\[
  h(t)=\FNV(\operatorname{id}(t)).
\]
Fix a length $N\in\{0,\dots,8\}$ (default $0$, i.e.\ off).  A \emph{T-suffix}
is a tuple $\TS=(H_1,\dots,H_N)\in(\B)^{N}$.  For a trace
$s_1,\dots,s_m$ of step hashes, define
\[
  H_k^{(0)}=\kappa_0,
  \qquad
  H_k^{(m)}=\mix\bigl(H_{k-1}^{(m-1)},s_m\bigr),
  \quad k=1,\dots,N,
  \label{eq:tsuffix}
\]
where $\kappa_0=\mathtt{0x5EED5EED}$ is the seed, $H_0^{(m-1)}$ is interpreted
as $\kappa_0$ for $k=1$, and
\[
  \mix(c,s)=
  \operatorname{avalanche}\bigl((c\oplus s)\cdot\mathtt{0x9E3779B9}\bigr),
\]
with the Murmur-style finalizer
\[
  \begin{aligned}
    h_1 &= \bigl(h\oplus(h\gg15)\bigr)\cdot\mathtt{0x85EBCA6B},\\
    h_2 &= \bigl(h_1\oplus(h_1\gg13)\bigr)\cdot\mathtt{0xC2B2AE35},\\
    \operatorname{avalanche}(h) &= h_2\oplus(h_2\gg16).
  \end{aligned}
\]
All shifts are logical right shifts on 32-bit words and all additions and
multiplications are modulo $2^{32}$.
\end{definition}

Written out step by step, the finalizer is: xor with $h\gg15$, multiply by
\texttt{0x85EBCA6B}, xor with $h\gg13$, multiply by \texttt{0xC2B2AE35}, xor
with $h\gg16$.  The recurrence gives the intended suffix semantics:
$H_1^{(m)}$ depends only on the last step, $H_2^{(m)}$ on the last two, and
$H_N^{(m)}$ on the last $N$ steps.  A trace shorter than $k$ contributes the
seed for the missing deeper history.  One arriving step updates all $N$
elements in $O(N)$ by a rolling carry, and a candidate edge is priced by
extending a copy of the current suffix with the edge's step.

\begin{definition}[Pair channels]\label{def:pairs}
For $N>0$ the colony keeps a second store
\[
  \Pair_a,\Pair_r\colon S\times(\B)^{N}\to[0,\Tau].
\]
A pair deposit at arrival marking $s$ with suffix $\TS$ lands flat, with no
back-propagation:
\[
  \Pair_c(s,\TS)\leftarrow
  \min\bigl(\Tau,\Pair_c(s,\TS)+A\bigr).
\]
Real-run milestones and catastrophes deposit the same amounts as their edge
deposits.  The pair channels evaporate with the same rates as the edge
channels, and an entry whose two channels both fall below $10^{-4}$ is
pruned.  With $N=0$ every pair operation is a no-op.
\end{definition}

The effect on the choice rule is exactly the $a^{+}$ and $r^{+}$ terms in
Equation~\eqref{eq:weight}: the pair attraction joins the base of the power,
the pair repulsion joins the exponent of the decay.  A catastrophe suffered
through one route therefore no longer condemns the same marking reached
through another route.

\subsection{One iteration}\label{sec:iteration}

\begin{algorithm}[H]
\footnotesize
\caption{One colony iteration (defaults in brackets)}
\label{alg:iteration}
\begin{algorithmic}[1]
\Require explored graph $G$, coverage database, start markings $S_{\mathrm{start}}$
\State refresh the off-target region if the graph grew \Comment{Definition~\ref{def:offtarget}}
\State refresh the target board and the predecessor index
\State $B_{\mathrm{exp}}\gets 200$ \Comment{expansion budget for this iteration}
\For{$i\gets1$ \textbf{to} $20$}
  \State pick a start marking uniformly from $S_{\mathrm{start}}$
  \State $\TS\gets(\kappa_0,\dots,\kappa_0)$; $s\gets$ start
  \For{$\mathit{step}\gets1$ \textbf{to} $40$}
    \If{$s$ is not expanded}
      \If{$B_{\mathrm{exp}}=0$} \textbf{break} \EndIf
      \State expand $s$; $B_{\mathrm{exp}}\gets B_{\mathrm{exp}}-1$
    \EndIf
    \State choose an edge $e$ by the two-step rule of Section~\ref{sec:choice}
    \If{no edge is enabled} \textbf{break} \EndIf
    \State $s\gets\tau(e)$; $\TS\gets\TS\cdot h(t_e)$
    \State apply the arrival rule of Section~\ref{sec:deposits}
    \If{the arrival rule stopped the ant} \textbf{break} \EndIf
  \EndFor
\EndFor
\If{evaporation is on}
  \State apply Equation~\eqref{eq:evaporation} and wipe saturated attraction
\EndIf
\end{algorithmic}
\end{algorithm}

The iteration is deterministic for a given seed: the colony owns a uniform
generator on top of \texttt{std::mt19937\_64}, so the result does not depend
on the standard library's distribution implementation.  A round may expand at
most 200 new markings in total, shared by all its ants.  The default number of
ants is 20 and the default walk length is 40 steps.

\section{The suggester}\label{sec:suggester}

The colony explores the model; the suggester chooses the next input for the
live SUT.  It scores each candidate controllable transition by a shortest
distance to a goal, corrected for risk, hazards and pheromone.  The two share
the graph, the coverage database and the target board.

\subsection{Goal distance}

A backward Dijkstra pass \cite{dijkstra1959note} from all goal markings over
the explored graph gives the cost
\[
  \dist(s)=\text{shortest number of edges from }s\text{ to a goal},
\]
with a tier penalty added for the goal's tier, so a slightly longer route to a
rare goal outranks a shorter route to an ordinary edge.  Goals are ranked in
tiers: rare transition, target, uncovered transition, uncovered edge,
frontier; $\dist(s)=\infty$ when no goal is reachable.  For a candidate $t$, let
$E_{\mathrm{C}}(t)$ be the set of edges in the current belief that fire $t$.
The best score over those edges is
\[
  \mathit{best}(t)=
  \min\Bigl(
    \min_{\substack{e\in E_{\mathrm{C}}(t)\\ d(e)=0}}
      \bigl(1+\tierpen(\operatorname{tier}(e))\bigr),\;
    \min_{e=(s,t,s')\in E_{\mathrm{C}}(t)}
      \bigl(1+\dist(s')\bigr)
  \Bigr),
\]
where the first minimum covers an uncovered edge at the belief itself and the
second covers a goal reached through the candidate.

\subsection{Robustness and hazard}

Let $B$ be the current belief set and let $E_{\mathrm{C}}(t)$ be the enabling
edges of $t$ in $B$.  The \emph{enabled fraction} is
\[
  \rho_{\mathrm{en}}(t)=\frac{|\{s\in B:t\text{ is enabled in }s\}|}{|B|}.
\]
For each enabling edge $e=(s,t,s')$, look at the environment moves out of
$s$; with $\omega$ the offered/taken weight and
$\mathbb{1}_{t}(s'')$ the indicator that $t$ is enabled after the environment
move to $s''$, define
\[
  \sigma(e)=
  \frac{\sum_{e'\in E_{\mathrm{O}}(s)}\omega(e')\,
        \mathbb{1}_{t}(\tau(e'))}
       {\sum_{e'\in E_{\mathrm{O}}(s)}\omega(e')},
\]
with $\sigma(e)=1$ when the denominator is zero.  The \emph{stability} is the
mean of $\sigma$ over $E_{\mathrm{C}}(t)$, and the \emph{robustness} is
\[
  \mathrm{rob}(t)=\rho_{\mathrm{en}}(t)\cdot
  \frac{1}{|E_{\mathrm{C}}(t)|}\sum_{e\in E_{\mathrm{C}}(t)}\sigma(e).
\]
The \emph{hazard} is the fraction of enabling edges whose target is
forbidden:
\[
  \mathrm{haz}(t)=
  \frac{|\{e\in E_{\mathrm{C}}(t):\cls(\tau(e))=\mathsf{Forbidden}\}|}
       {|E_{\mathrm{C}}(t)|}.
\]

\subsection{The effective score}

With the pheromone sums
\[
  A(t)=\sum_{e\in E_{\mathrm{C}}(t)}a(e),\quad
  R(t)=\sum_{e\in E_{\mathrm{C}}(t)}r(e),\quad
  G(t)=\sum_{e\in E_{\mathrm{C}}(t)}g(e),
\]
and the corresponding pair sums $A^{+}(t),R^{+}(t)$ taken over the arrival
contexts of the edges, the suggester's \emph{effective score} is
\[
  \mathrm{eff}(t)=
  \mathit{best}(t)
  + \rho_{\mathrm{risk}}\bigl(1-\mathrm{rob}(t)\bigr)
  + h_{\mathrm{haz}}\,\mathrm{haz}(t)
  + w_{\mathrm{rep}}\frac{R(t)+G(t)+R^{+}(t)}{|E_{\mathrm{C}}(t)|\Tau}
  - w_{\mathrm{att}}\frac{A(t)+A^{+}(t)}{|E_{\mathrm{C}}(t)|\Tau},
  \label{eq:suggester}
\]
with defaults $\rho_{\mathrm{risk}}=2.0$, $h_{\mathrm{haz}}=20.0$,
$w_{\mathrm{rep}}=3.0$ and $w_{\mathrm{att}}=1.0$.  The candidate with the
smallest effective score is suggested; ties are broken by higher robustness,
then by a deterministic transition order (or a seeded random choice when a
tie-break seed is configured).  The pheromone terms are scaled by
$|E_{\mathrm{C}}(t)|\Tau$, so a full-strength channel on every enabling edge
contributes exactly the configured weight in steps.

\begin{remark}
Equation~\eqref{eq:suggester} is the coupling between the two halves of the
heuristic.  The ants learn on the model; the suggester spends that knowledge
on the next real input.  The static off-target channel $G(t)$ enters with the
repellent pheromone, so an input whose every belief edge leads into the
off-target region is as unattractive as one that leads into a catastrophe.  A
hard veto can still refuse a candidate whose every effect is forbidden; the
phase policy \texttt{veto} is the default.
\end{remark}

\section{Abstract algorithms}\label{sec:algorithms}
% DeepSeek V4.1 Flash: abstract algorithms requested by Egge (journal §38).

The definitions above are declarative.  This section restates them as the four
procedures a reader needs in order to reproduce a campaign: the lazy
construction of the reachability graph, the walk of one ant, the two pheromone
updates, and the selection of the next event for the live SUT.  The
pseudo-code stays close to the implementation but abstracts from storage: a
marking is an interned node, and $\mathrm{id}(M)$ denotes the node that
carries $M$.  Algorithm~\ref{alg:iteration} gave the outer colony iteration;
the procedures below give its inner steps.

\subsection{Reachability-graph generation}\label{sec:alg-rg}

Algorithm~\ref{alg:rg} builds the explored subgraph of
Definition~\ref{def:explored} on demand.  A node is created when a firing
first reaches its marking; an existing marking is reused, so two paths that
meet share one node.  The outer loop takes unexpanded nodes in creation order,
which is breadth first because identifiers are handed out in discovery order.

\begin{algorithm}[H]
\footnotesize
\caption{Lazy generation of the explored reachability graph}
\label{alg:rg}
\begin{algorithmic}[1]
\Require net $N=(P,T,F,I,w,\beta,\ell,\kappa)$, initial markings
         $\mathcal{M}_0$, limit $B$
\Ensure explored subgraph
        $(S_{\mathrm{known}},S_{\mathrm{exp}},E_{\mathrm{exp}})$ and frontier
        $U=S_{\mathrm{known}}\setminus S_{\mathrm{exp}}$
\State $S_{\mathrm{known}}\gets\emptyset$;
       $S_{\mathrm{exp}}\gets\emptyset$;
       $E_{\mathrm{exp}}\gets\emptyset$;
       $S_0\gets\emptyset$
\ForAll{$M\in\mathcal{M}_0$}
  \State insert $M$ into $S_{\mathrm{known}}$ and $S_0$ unconditionally
\EndFor
\While{$\exists s\in S_{\mathrm{known}}\setminus S_{\mathrm{exp}}$}
  \State take the unexpanded $s$ with the smallest identifier
  \State \Call{Expand}{$s$}
\EndWhile
\State \Return $(S_{\mathrm{known}},S_{\mathrm{exp}},E_{\mathrm{exp}})$
\Statex
\Procedure{Expand}{$s$}
  \State $E_s\gets\emptyset$; $n_{\mathrm{en}}\gets0$;
         $\mathit{cut}\gets\mathsf{false}$
  \ForAll{$t\in T$}
    \If{$\neg\Enabled(t,s)$} \textbf{continue} \EndIf
    \State $n_{\mathrm{en}}\gets n_{\mathrm{en}}+1$
    \If{$\mathsf{lifecycle}(t)\ne\bot\wedge\post_t=\emptyset$}
      \State \textbf{continue} \Comment{engine-landed, no token-game edge}
    \EndIf
    \State $(\mathit{status},M')\gets\Call{Fire}{s,t}$
    \If{$\mathit{status}\ne\mathsf{Ok}$}
      \State record the bound violation; \textbf{continue}
    \EndIf
    \If{$\mathrm{id}(M')=\bot\wedge|S_{\mathrm{known}}|\ge B$}
      \State $\mathit{cut}\gets\mathsf{true}$; \textbf{continue}
    \EndIf
    \If{$\mathrm{id}(M')=\bot$}
      \State create and classify the node $M'$ in $S_{\mathrm{known}}$
    \EndIf
    \State append $e=(s,t,M')$ to $E_s$
  \EndFor
  \State add $s$ to $S_{\mathrm{exp}}$ with enabled count $n_{\mathrm{en}}$
         and flag $\mathit{cut}$
  \State $E_{\mathrm{exp}}\gets E_{\mathrm{exp}}\cup E_s$
\EndProcedure
\end{algorithmic}
\end{algorithm}

Initial markings are inserted even when the limit has already been reached;
the cap governs only nodes created by firing.  A firing that would exceed a
place bound, or overflow an unbounded place, is recorded as a bound violation
and produces no edge.  A lifecycle input with no output arcs counts towards
$n_{\mathrm{en}}$---the marking is not a dead end, because the engine can
restart the SUT---but creates no token-game edge.  When the cap prevents a
successor from being created, the source is flagged incomplete and the graph
is marked truncated, so analyses cannot mistake the frontier for a deadlock.

\subsection{Ant movement}\label{sec:alg-ant}

Algorithm~\ref{alg:ant} follows one ant from a start marking.  The walk stops
when the budget or the step limit is exhausted, no edge is enabled, or the
arrival rule ends the ant.  \textproc{Choose} is the two-step draw of
Section~\ref{sec:choice}; Algorithm~\ref{alg:arrive} is the arrival rule of
Section~\ref{sec:deposits}.  The pair channels are zero when the T-suffix
layer is off.  In \textproc{Choose}, $U$ is a fresh uniform draw on $[0,1)$.

\begin{algorithm}[H]
\footnotesize
\caption{One ant: walk and choice}
\label{alg:ant}
\begin{algorithmic}[1]
\Procedure{Walk}{$s_0$}
  \State $s\gets s_0$; $\TS\gets(\kappa_0,\dots,\kappa_0)$
  \For{$k\gets1$ \textbf{to} $40$}
    \If{$s\notin S_{\mathrm{exp}}$}
      \If{$B_{\mathrm{exp}}=0$}
        \State \Return exhausted
      \EndIf
      \State \Call{Expand}{$s$}; $B_{\mathrm{exp}}\gets B_{\mathrm{exp}}-1$
    \EndIf
    \State $e\gets\Call{Choose}{s,\TS}$
    \If{$e=\bot$}
      \State \Return exhausted
    \EndIf
    \State $s\gets\tau(e)$; $\TS\gets\TS\cdot h(t_e)$
    \If{\Call{Arrive}{$s$}}
      \State \Return ended
    \EndIf
  \EndFor
  \State \Return exhausted
\EndProcedure
\Statex
\Function{Choose}{$s,\TS$}
  \State $I\gets\{e:\sigma(e)=s,\ \kappa(t_e)=\mathsf{C}\}$
  \State $O\gets\{e:\sigma(e)=s,\ \kappa(t_e)\ne\mathsf{C}\}$
  \If{$I=O=\emptyset$}
    \State \Return $\bot$
  \EndIf
  \If{$I=\emptyset$}
    \State $P\gets O$
  \ElsIf{$O=\emptyset$}
    \State $P\gets I$
  \ElsIf{$U<\rho_{\mathrm{env}}$}
    \State $P\gets O$
  \Else
    \State $P\gets I$
  \EndIf
  \ForAll{$e\in P$}
    \If{$e\in O$}
      \State $\wgt(e)\gets\omega(e)$
    \Else
      \State $\TS'\gets\TS\cdot h(t_e)$
      \State $a^+\gets\Pair_a(\tau(e),\TS')$;
             $r^+\gets\Pair_r(\tau(e),\TS')$
      \State $\wgt(e)\gets
             (a(e)+a^++0.1)^{\alpha}\nu(e)^{\beta}
             \exp[-\lambda(r(e)+g(e)+r^+)]$
    \EndIf
  \EndFor
  \State draw $e\in P$ with probability
         $\wgt(e)/\sum_{e'\in P}\wgt(e')$; if the sum is zero, uniformly
  \State \Return $e$
\EndFunction
\end{algorithmic}
\end{algorithm}

\begin{algorithm}[H]
\footnotesize
\caption{Arrival rule of one ant}
\label{alg:arrive}
\begin{algorithmic}[1]
\Function{Arrive}{$s$}
  \If{$\cls(s)=\mathsf{Forbidden}$}
    \State \Call{Deposit}{$r,s,c_{\mathrm{cat}}$};
           \Return $\mathsf{true}$
  \ElsIf{$\cls(s)=\mathsf{Seeking}$}
    \State \Call{Deposit}{$a,s,c_{\mathrm{tgt}}$};
           \Return $\mathsf{true}$
  \EndIf
  \If{$s\notin S_{\mathrm{exp}}$}
    \If{$B_{\mathrm{exp}}>0$}
      \State \Call{Expand}{$s$}; $B_{\mathrm{exp}}\gets B_{\mathrm{exp}}-1$
    \Else
      \State \Call{Deposit}{$a,s,c_{\mathrm{unexp}}$};
             \Return $\mathsf{true}$
    \EndIf
  \EndIf
  \ForAll{$e\in E_{\mathrm{exp}}$ with $\sigma(e)=s$}
    \If{$\cls(\tau(e))\ne\mathsf{Forbidden}\wedge d(e)=0$}
      \State \Call{Deposit}{$a,s,c_{\mathrm{unexp}}$};
             \Return $\mathsf{true}$
    \EndIf
  \EndFor
  \State \Return $\mathsf{false}$
\EndFunction
\end{algorithmic}
\end{algorithm}

The default amounts are $c_{\mathrm{unexp}}=1$, $c_{\mathrm{tgt}}=2$ and
$c_{\mathrm{cat}}=3$.  The expansion budget is shared by all ants of an
iteration.  A forbidden arrival deposits into the repel channel and ends the
ant; a seeking target deposits into the attract channel and ends it.  An
unexpanded marking is expanded while budget remains; with no budget left it
is deposited as unexplored.  A marking with an uncovered outgoing edge is
also a reason to stop.  The static off-target region does not kill the ant:
behaviour inside it may still be worth covering.

\subsection{Pheromone setting}\label{sec:alg-pheromone}

Two operations write the learned channels; the static off-target channel is
not learned.  \textproc{Deposit} walks the predecessor index backwards from the
arrival marking for $D$ levels (default $D=5$) and adds $A\gamma^d$ to every
reverse edge at depth $d=0,\dots,D-1$, capped at $\Tau$.  A pair deposit is
flat.  \textproc{Evaporate} applies one evaporation step and wipes the attraction
of saturated edges; with the matured-target exception, the way into a
$\cls(\mathsf{Forbidden})$ marking is spared.

\begin{algorithm}[H]
\footnotesize
\caption{Pheromone deposition, back-propagation and evaporation}
\label{alg:pheromone}
\begin{algorithmic}[1]
\Procedure{Deposit}{$c,s,A$}
  \State $Q\gets\{e\in E_{\mathrm{exp}}:\tau(e)=s\}$
  \For{$d\gets0$ \textbf{to} $D-1$}
    \ForAll{$e\in Q$}
      \State $c(e)\gets\min(\Tau,c(e)+A\gamma^{d})$
    \EndFor
    \State $Q\gets$ edges entering a source not visited at a
           smaller depth
  \EndFor
\EndProcedure
\Statex
\Procedure{PairDeposit}{$c,s,\TS,A$}
  \State $\Pair_c(s,\TS)\gets\min(\Tau,\Pair_c(s,\TS)+A)$
\EndProcedure
\Statex
\Procedure{Evaporate}{$\mathit{spare}$}
  \ForAll{$e\in E_{\mathrm{exp}}$}
    \If{$\mathrm{saturated}(e)$}
      \State $a(e)\gets0$
    \EndIf
    \If{$\neg\mathit{spare}(\tau(e))$}
      \State $a(e)\gets(1-\epsilon_a)a(e)$;
             $r(e)\gets(1-\epsilon_r)r(e)$
    \EndIf
  \EndFor
  \ForAll{$(s,\TS)$ in the pair store}
    \If{$\neg\mathit{spare}(s)$}
      \State age both pair channels by the same factors
    \EndIf
    \State delete the key if both channels are below $10^{-4}$
  \EndFor
\EndProcedure
\Statex
\Procedure{RefreshRegion}{}
  \State $S_{\mathrm{reach}}\gets$ markings that can reach a seed of
         Definition~\ref{def:offtarget}
  \State $g(e)\gets\rho$ if
         $\tau(e)\in S_{\mathrm{exp}}\setminus S_{\mathrm{reach}}$, else $0$
\EndProcedure
\end{algorithmic}
\end{algorithm}

The direct way into an event carries the full amount; a predecessor five
steps behind carries $\gamma^5\approx0.168$ of it.  The default rates are
$\epsilon_a=0.10$ and $\epsilon_r=0.02$, so a catastrophe is remembered longer
than a success.  Real-run milestones and catastrophes use the same
\textproc{Deposit}, with the amounts listed in Section~\ref{sec:deposits}, and
also write the pair store.  The off-target value
$\rho$ is set by \textproc{RefreshRegion} whenever the graph grew; it is never
accumulated and never evaporated.

\subsection{Event selection in the tracker}\label{sec:alg-event}

The belief tracker does not choose the next event by itself.  It maintains
$B$, closes it under silent steps, and lists the controllable transitions
enabled in at least one marking of $B$; the suggester turns that list into an
event.  Algorithm~\ref{alg:event} is the selection rule, with the colony
optional.  The lifecycle gate is checked first: while the SUT is paused or
down, only the lifecycle input that brings it back is returned.  The optional
rare hunt is a second early branch: if an unseen rare non-controllable
transition is enabled in some belief marking, the suggester returns a dwell
(or a wait when its wait-for is zero) instead of an input; a controllable
rare input is handled by the ordinary score and sent after its dwell.

\begin{algorithm}[H]
\footnotesize
\caption{Selecting the next event from the tracker's belief}
\label{alg:event}
\begin{algorithmic}[1]
\Require belief $B$; candidates $\mathcal{C}$ from the tracker; board
         $\cls$; colony $A$ or $\bot$; live trace $\TS$
\Ensure a suggestion, or $\mathsf{Wait}$, $\mathsf{Done}$ or
        $\mathsf{Stuck}$
\If{the lifecycle gate is not $\mathsf{Running}$}
  \State \Return the lifecycle input the gate allows, or $\mathsf{Stuck}$
\EndIf
\If{the rare hunt is on and an unseen rare non-controllable transition
    with attempts left is enabled in some policy-allowed $s\in B$}
  \State \Return $\mathsf{Dwell}$ (or $\mathsf{Wait}$ when the wait is zero)
\EndIf
\ForAll{$t\in\mathcal{C}$}
  \State $E_{\mathrm{C}}(t)\gets
         \{(s,e):s\in B,\ e=\text{the edge of }t\text{ at }s\}$
\EndFor
\State expand up to $B_{\mathrm{exp}}$ unexpanded successors, starting from the
       candidate targets
\State refresh the board and the predecessor index
\State seed the goal costs and tiers, then relax backwards
       \Comment{Section~\ref{sec:suggester}}
\ForAll{$t$ with $E_{\mathrm{C}}(t)\ne\emptyset$}
  \State compute $\mathrm{rob}(t)$ and $\mathrm{haz}(t)$
         \Comment{Section~\ref{sec:suggester}}
  \If{the policy is $\mathsf{Veto}$ and every target is forbidden}
    \State veto $t$; \textbf{continue}
  \EndIf
  \State $\mathit{best}(t)\gets$ best score over $E_{\mathrm{C}}(t)$
         \Comment{tier penalty included, Section~\ref{sec:suggester}}
  \State $\mathrm{goal}(t)\gets$ tier of that best score
  \State $\mathrm{eff}(t)\gets\mathit{best}(t)
         +\rho_{\mathrm{risk}}(1-\mathrm{rob}(t))
         +h_{\mathrm{haz}}\mathrm{haz}(t)$
  \If{$A\ne\bot$ and $\Tau>0$}
    \State $\mathrm{eff}(t)\gets\mathrm{eff}(t)
           +w_{\mathrm{rep}}\frac{R(t)+G(t)+R^+(t)}
             {|E_{\mathrm{C}}(t)|\Tau}
           -w_{\mathrm{att}}\frac{A(t)+A^+(t)}
             {|E_{\mathrm{C}}(t)|\Tau}$
  \EndIf
\EndFor
\If{no candidate has a finite score}
  \State \Return $\mathsf{Wait}$, $\mathsf{Done}$ or $\mathsf{Stuck}$
         \Comment{by the goal and expected-event test}
\EndIf
\If{the rare hunt is on and some $t$ has
    $\mathrm{goal}(t)=\mathsf{Rare}$}
  \State keep only the candidates with $\mathrm{goal}(t)=\mathsf{Rare}$
\EndIf
\State choose the least $\mathrm{eff}$; break ties by robustness, goal tier,
       transition order
\If{the chosen candidate is rare with a positive wait-for}
  \State \Return $\mathsf{Dwell}$ for the chosen input
\EndIf
\State \Return the chosen input
\end{algorithmic}
\end{algorithm}

A frontier marking is seeded at $\tierpen(\mathsf{Frontier})$, a seeking
marking at $\tierpen(\mathsf{Target})$, and every source of an uncovered edge
at $1+\tierpen(\operatorname{tier}(e))$.  The backward Dijkstra over the
predecessor index then fills $\dist$ for the remaining markings.

\paragraph{With and without pheromones.}
When the campaign runs with the colony (the default), the session performs
\texttt{antIterationsPerSuggestion} iterations (default one) on the current
belief before Algorithm~\ref{alg:event}, and the attached colony contributes
$A(t),R(t),G(t)$ and, with the T-suffix layer on, $A^+(t),R^+(t)$.  With
\texttt{--no-ants}, no ants walk and no colony is attached: the conditional
line in Algorithm~\ref{alg:event} is omitted and the ranking is the baseline
\[
  \mathrm{eff}_{\mathrm{base}}(t)=\mathit{best}(t)
  +\rho_{\mathrm{risk}}\bigl(1-\mathrm{rob}(t)\bigr)
  +h_{\mathrm{haz}}\,\mathrm{haz}(t).
\]
Everything else---belief, candidates, goal distance, robustness, hazard and
the tie-break---is identical.  The colony therefore changes which input is
chosen, not what counts as a legal or a useful one.

\section{Application: what the campaigns found}\label{sec:application}

The definitions above were not written after the fact.  They describe a tool
that has been run on a ladder of models, from a four-place drink machine to a
ten-thousand-state product of four communicating machines, and the engineering
record in \texttt{journal.md} contains the measurements.  This section revisits
that record.  It is an application chapter, not a controlled evaluation: the
models were built to exercise one mechanism at a time, the campaigns are
deterministic only for a given seed, and the wall-clock numbers belong to the
development machine.  The purpose is to show which parts of the machinery were
actually exercised and what the numbers do and do not say.

\subsection{A ladder of models}

The examples shipped with the tool form a deliberate ladder.

\begin{itemize}[leftmargin=2.2em]
  \item \texttt{examples/drink-machine.json} is the smallest complete model:
        a coin, a drink, a return, and one target marking.  It is the model
        used in the first end-to-end campaign and in the reference agent.
        The same net is shipped as \texttt{examples/coffee-pnml.pnml}, so the
        PNML reader can be checked against a JSON model whose behaviour is
        known.
  \item \texttt{examples/elevator-door.json} adds place bounds, inhibitor
        arcs and a \texttt{covers}-matched error marking (\texttt{jammed})
        whose post-mortem is a snapshot.  It is the smallest model on which
        the belief tracker must distinguish a forbidden marking from an
        ordinary dead end.
  \item \texttt{examples/atm.json} is the worked example that exercises the
        whole toolbox: 22 places, 22 transitions, three initial markings,
        two target markings, an error marking, four error transitions, a
        traced internal announcement, an unpredictable human step and a
        \texttt{covers} target that requires a real withdrawal.  It is
        discussed below.
  \item the four-machine ring in \texttt{tests/tst\_comm10k.cpp} is the scale
        model: four machines sharing four binary channels, exactly
        $5^4\cdot2^4=10{,}000$ markings and $64{,}000$ edges, with one goal
        and one error marking per machine.  It is used for the run-length
        experiment, not for a functional story.
\end{itemize}

The ladder matters because each model pins a different claim.  A small model
can make a campaign look better than it is; a large model can hide a
mechanism in noise.  The journal records both, and the publication quotes the
small ones only for structure.

\subsection{The ATM campaign}\label{sec:atm}

The ATM model is the closest the project comes to an application.  A visit
inserts a card, reads it, asks for a PIN, optionally opens a bank session and
a database connection, checks the card against a stolen report, answers a
withdrawal request, and ends with cash, a rejection, a confiscated card or a
component failure.  The model has three initial markings,
\texttt{atm-lobby}, \texttt{atm-drivein} and \texttt{atm-branch}.  They differ
in the number of tokens in the cash cassette (2, 1 and 0), because identical
markings would be interned into one graph node and the round-robin over the
three ATMs would be invisible in the run log.  With the cassette levels the
rotation is observable, and the branch machine, which is out of cash, cannot
reach the good target at all.

The target \texttt{visit-done} is matched with \texttt{covers} and requires
both \texttt{visit\_done} and \texttt{cash\_taken}; the plain card-return
exit reaches \texttt{visit-aborted}, which is a target marking but not an
intentional one, so it does not end the campaign.  The error marking
\texttt{card-swallowed} is reached after three invalid PINs, is matched with
\texttt{covers}, has policy \texttt{until\_covered} and one crash dump.  Four
error transitions---\texttt{confiscate}, \texttt{session\_fail},
\texttt{db\_fail} and \texttt{atm\_fail}---sit where the belief is a single
state and their competitors are visible, so a certain step can be attributed
and dumped.  The PIN input is tagged \texttt{unpredictable}, which keeps the
catastrophe deposit flat on a way in that measured a human's reaction time.

A campaign with \texttt{--sim-seed 5 --max-runs 45} tells the deterministic
story recorded in journal \S13: 37 runs, five crash post-mortems (the four
error transitions, each dumped once, and the swallowed-card marking), the
error marking hit twice and satisfied, the good target reached at run 37, and
the three initial markings used 13, 12 and 12 times.  Other simulator seeds
produce the same story in 37--58 runs.  The point is not the exact count.  The
point is that one model exercised initial-marking rotation, target and error
policies, certain error transitions, post-mortem budgets, the
\texttt{covers} match, the unpredictable tag and the run log's trace view
without a special-case code path for the ATM.

The model also shows the price of ambiguity.  The two silent routes into the
reading state are indistinguishable from outside, so \texttt{reader\_check}
and \texttt{pad\_check} can never be covered \emph{definitely}; they are
credited only when every consistent path crosses them.  This is not a bug in
the colony.  It is the belief tracker being honest about what the SUT let it
observe, and it is the reason the model makes the stolen-card announcement a
traced internal step: the announcement narrows the belief before the
observable \texttt{confiscate} step, so that error transition is caught
certainly.

\subsection{The fork benchmark: what one edge channel cannot learn}

The central negative result is the fork benchmark of journal \S25.  A short
corridor leads to a fork.  One branch is a single input, \texttt{mine}, that
crashes the SUT; the other is a two-step safe route to the same region.  The
region is a ring of 1000 markings with a skip edge, 1008 markings and 2010
edges in all.  Every run starts in the corridor and meets the same fork with
the same pheromone memory.  The crash is a predicted target while it is still
being reproduced, so the target board deliberately keeps the way in
attractive; after the experiment resolves, the catastrophe channel takes
over.  The campaign budget is six runs of 200 events.

The measured comparison (three seeds, graph seed 7) is stark:

\begin{center}
\small
\begin{tabular}{@{}lrrr@{}}
\toprule
guidance & crashes & long runs & covered edges \\
\midrule
no ants & 5 & 1 & 201 \\
ants, edge channel only & 5 & 1 & 201 \\
ants, edge and T-suffix & 3 & 3 & 585--586 \\
\bottomrule
\end{tabular}
\end{center}

Without the colony, the campaign explores the region once and then walks into
the crashing shortcut at the start of each following run.  The edge channel
alone reproduces that outcome exactly, on every seed.  The reason is
saturation: the target walks keep passing through the region's entry state on
their safe way to the target, so the shared per-edge values around the fork
sit at the cap $\Tau=10$ on both sides.  The later catastrophe deposit lands
on saturated values and cannot make the shortcut look worse than its
neighbours.  The pair layer does not share that fate.  Its key is the marking
together with the recent trace suffix, so the crash's pair is ``entry state,
suffix ending in \texttt{mine}'', while a safe arrival lands on a different
pair.  The pair keeps its catastrophe repel while the walks flood the
neighbourhood.  The result is about 2.9 times the coverage for the same run
budget, at about 2.6 times the wall-clock time.

The lesson is not ``ants help''.  It is that one value per state cannot carry
``this way in is fatal'' while the target walks keep saying ``this state is on
the way''.  The second memory, keyed by how the state was reached, can.  This
is exactly the role of the T-suffix layer in Section~\ref{sec:tsuffix}, and it
is why the layer is optional: on models where the safe and the bad route do
not share their recent history, the edge channel is enough, and the pair store
only costs memory.

\subsection{Run length and evaporation}

Two further experiments bound the tuning.  The first is the evaporation
boundary (journal \S27).  With per-suggestion aging, a run of 200 suggestions
keeps $0.98^{200}\approx1.8\%$ of the repel lesson its beginning laid; the
repel half-life is about 34 suggestions.  Moving the aging to the run
boundary---\texttt{--evaporate-between-runs}---keeps the lesson intact inside
a run but lets a wrong lesson steer the whole run.  On the fork benchmark the
boundary variant covered 777 edges instead of 586, with two crashes instead of
three and four long runs instead of three.  The same contrast held at 400
events.  The default remains per-suggestion aging, because a long run is
exactly where a stale lesson is most expensive, and because boundary aging
makes a campaign sensitive to the run length.

The second is the run-length experiment on the 10{,}000-state ring (journal
\S32).  Four campaigns with a ten-minute budget and \texttt{--max-events} 50,
200, 800 and 3200 covered 12{,}198, 12{,}527, 13{,}254 and 11{,}798 edges on
average.  The middle range was best on average, but the best-to-worst
difference was only 12\%, while the seed-to-seed spread within one size
reached 41\%.  The reliable part was the per-event rate: edges per event rose
from 0.918 at 50 events to 0.972 at 3200, because a short run re-tests the
initial region on every restart.  Run length is therefore a weak coverage
lever; it is a lever on how much of a bad run has to be repeated and how often
the SUT setup is paid.

\subsection{Rare transitions and the hunt}

The rare-transition machinery (journal \S30) is the one place where the model
declares a probability rather than measuring one.  A transition may carry
\texttt{rare} with a declared $p$, an attempt budget and an optional dwell.
The number $p$ is descriptive: the engine cannot observe the SUT's real odds,
so it never schedules from $p$.  The hunt is driven by the attempt budget.
With \texttt{--hunt-rare}, an unseen rare transition whose budget is not spent
becomes a goal above targets, fresh transitions, fresh edges and the frontier;
an equal-cost tie also goes to the rare goal, or a fresh edge on the last step
would keep the hunt from making the choice it was asked to make.  The enabling
state, not the edge, is the goal for a move of the SUT, because the tester
cannot take that edge.  A dwell at the enabling state suspends the idle
timeout, is ended early by any real event, and charges one attempt; a
controllable rare input with a non-zero dwell is withheld until the dwell
ends.  A \texttt{Quiescent} report is not an event and does not end the dwell.

The honest limits are part of the feature.  An unreachable rare transition
keeps the campaign alive until a hard limit, because no attempt can be spent
and the hunt is still pending.  Attempt budgets do not survive a run log, so
a continued campaign starts them fresh while the transition hits survive.
A rare move that the SUT makes only under conditions the model does not
describe cannot be hunted by waiting at a state the model believes in.  The
hunt is not a probability oracle.

\subsection{What the applications do not show}

The ladder is not a representative sample of industrial models.  The coffee
machine, the elevator door and the ATM are small enough to inspect by hand;
the ring is regular enough to count its state space exactly.  None of them has
data variables, real-time constraints, a distributed clock or a network
partition.  The tool deliberately leaves those to the command hooks and to the
SUT's own oracle (Section~\ref{sec:implementation}), and the journal's limits
chapter records where the heuristic stops helping: saturation of a sought
target, exact-match sparsity of the pair key, truncated graphs, human and
timing unpredictability, and a model that is simply wrong.

Two issues that the campaigns exposed are worth naming here because they are
easy to miss in a definition paper.  First, several initial markings are only
as useful as the plug-in's ability to realise them: the tracker sends the
marking's name in \texttt{start\_run}, and a wrapper that ignores it can start
the SUT in the wrong state while the belief is built from the model's marking.
The ATM's cassette levels make the rotation observable in the log, but they
cannot make a wrapper honest.  At version \texttt{v0.0.8} a plug-in can
advertise the names it can realise in its hello, and the tracker filters the
rotation and refuses an explicit \texttt{--initial} outside the list (journal
\S41); the capability is a promise, not a proof.  Second, PNML cannot restore
the testing annotations: it has one initial marking and no targets, errors,
milestones, bounds or crash predictions.  The tool accepts PNML for
\texttt{check} and \texttt{run}, and \texttt{eggembt import} now reads it
through the same validating loader and writes the eggembt JSON model that is
to be annotated, saying what the conversion cannot carry (journal \S41).
PNML's expressive poverty is a property of the format, not of the import; the
JSON model remains the complete specification.  Both issues are recorded in
\texttt{TODO.md} and in journal \S39.

\section{Parameters and correspondence to the implementation}\label{sec:implementation}

\subsection{Default parameters}

\begin{center}
\small
\begin{tabular}{@{}llr@{}}
\toprule
Symbol & Meaning & Default \\
\midrule
$n_{\mathrm{ants}}$ & ants per iteration & 20 \\
$\mathit{steps}$ & steps per ant & 40 \\
$B_{\mathrm{exp}}$ & markings expanded per iteration & 200 \\
$D$ & back-propagation depth & 5 \\
$\gamma$ & depth decay & 0.7 \\
$\alpha,\beta,\lambda$ & attract, novelty, repel exponents & 1, 1, 1 \\
$\varepsilon_a$ & attractive evaporation per iteration & 0.10 \\
$\varepsilon_r$ & repellent evaporation per iteration & 0.02 \\
$3.0$ & catastrophe deposit (ants) & 3.0 \\
$2.0$ & target deposit & 2.0 \\
$1.0$ & unexplored deposit & 1.0 \\
$1.5$ & milestone deposit per rank & 1.5 \\
$\rho$ & off-target repel & 1.0 \\
$\Tau$ & pheromone cap & 10.0 \\
$p_{\mathrm{env}}$ & environment-move probability & 0.5 \\
$N$ & T-suffix length & 0 (off), max 8 \\
$\theta$ & coverage threshold & 1 \\
$\rho_{\mathrm{risk}}$ & suggester risk weight & 2.0 \\
$h_{\mathrm{haz}}$ & suggester hazard weight & 20.0 \\
$w_{\mathrm{att}},w_{\mathrm{rep}}$ & suggester pheromone weights & 1.0, 3.0 \\
\bottomrule
\end{tabular}
\end{center}

\subsection{Symbols and code}

\begin{center}
\small
\begin{tabular}{@{}p{0.28\linewidth}p{0.62\linewidth}@{}}
\toprule
Mathematics & Implementation (\texttt{src/}) \\
\midrule
$N=(P,T,F,I,w,\beta,\ell,\kappa)$ & \texttt{model/Net.h}, \texttt{NetBuilder} \\
$M\fir{t}M'$ & \texttt{engine/Semantics.cpp}: \texttt{isEnabled}, \texttt{fire} \\
$\kind(M)$, $\models_{\mathsf{exact}}$, $\models_{\mathsf{covers}}$ &
\texttt{engine/Semantics.cpp}: \texttt{classify}, \texttt{matches} \\
$G_{\mathrm{exp}}$, frontier, incomplete &
\texttt{engine/ReachabilityGraph.h} \\
$d(e)$, $\theta$, $\nu(e)$ & \texttt{engine/CoverageDb.h},
\texttt{isSaturatedEdge} \\
$\cls(s)$, target phases & \texttt{engine/TargetBoard.h} \\
$\Off$, $g(e)$ & \texttt{engine/AntColony.cpp}:
\texttt{refreshOffTargetRegionIfGrown}, \texttt{region} \\
$a(e),r(e)$ & \texttt{engine/AntColony.h}: \texttt{attract}, \texttt{repel} \\
$\wgt(e)$, the two-step draw & \texttt{AntColony::weight}, \texttt{AntColony::choose} \\
deposits and back-propagation & \texttt{AntColony::deposit},
\texttt{ReverseIndex::backwardWalk} \\
T-suffix, $\mix$, pair store & \texttt{engine/TSuffix.h},
\texttt{AntColony::pairAttract} \\
$\mathrm{eff}(t)$ & \texttt{engine/Suggester.cpp}, \texttt{Suggester::suggest} \\
\bottomrule
\end{tabular}
\end{center}

\subsection{What is not in this document}

The document defines the model and the exploration heuristic, not the testing
theory around them.  The plug-in wire protocol, the run log format, the
post-mortem machinery, the crash-credibility experiment, the temporal-logic
open item and the command hooks are described in \texttt{journal.md} and
\texttt{docs/PLUGIN\_PROTOCOL.md}.  The publication is intentionally
implementation-faithful: the defaults and the edge cases quoted here are the
ones in the source at \texttt{v0.0.8}.  If a definition here and the code ever
disagree, the code is the ground truth and this document should be corrected.

\section{Conclusion}

We have given a compact mathematical account of the two layers that make
\texttt{eggembt} more than a random walk: a typed, weighted, bounded Petri net
with target annotations and a lazily explored reachability graph; and an ant
colony whose edge weights combine attractive pheromone, repellent pheromone, a
static off-target region, novelty and an optional trace-suffix memory.  The
off-target region is the set of expanded markings from which no target,
sought fault or unknown territory is reachable; it is reward-inaccessible and
closed under successors, and its incoming edges carry a fixed repulsion.  The suggester then turns the colony's
learned channels into a one-step score for the live SUT.  The definitions are
deliberately small enough to check against the implementation, and every
symbol has a named counterpart in the source.

\bibliographystyle{plain}
\bibliography{references}

\end{document}
